\chapter{INTRODUCTION}
\justifying

Large-scale distributed architectures and data-driven security systems are becoming more and more common features of modern computing environments. Ensuring system reliability and data confidentiality has become a crucial challenge as software infrastructures become more complex and sensitive information is shared across digital platforms. Machine learning-driven automated analysis methods are now essential for tackling these issues, especially in the areas of system monitoring and cryptographic security assessment. Log-based anomaly detection in distributed systems, vulnerability analysis of secret image sharing schemes under learning-based attacks, and cryptanalysis of chaos-based medical image encryption schemes are three complementary security problems that are integrated in this project.

Finding anomalies in system logs is the focus of the first part. Runtime events, state transitions, and error conditions that represent system behavior are documented in system logs. Conventional parsing-based log analysis techniques are format-sensitive and rely on structured templates. Although parsing-free methods increase robustness, they frequently treat all tokens equally and ignore numerical data, which results in the loss of quantitative and structural context. In order to create structured message representations while maintaining numeric tokens, this project suggests a role-aware representation learning framework that directly induces latent token roles from anomaly supervision. Improving detection accuracy without using explicit log parsing is the goal.

The second part assesses how resilient XOR-based secret image sharing schemes are to contemporary machine learning threats. Information-theoretic guarantees that no information about the secret image can be revealed by any strict subset of shares are provided by classical $(n, n)$ XOR schemes. These assurances, however, are based on conventional threat models that do not take data-driven reconstruction methods into account. This project investigates whether an adversary based on autoencoders can use only a subset of a secret image's shares to reconstruct meaningful approximations of the image. To determine whether there is practical information leakage, reconstruction performance is evaluated using common image similarity metrics.

The third part focuses on the cryptanalysis of a chaos-based medical image encryption scheme. Recent encryption methods designed for medical images often rely on chaotic maps combined with permutation-diffusion structures to achieve confusion and diffusion. While such schemes demonstrate strong statistical properties, they frequently lack rigorous resistance against cryptographic attacks. This project analyzes a medical image encryption algorithm that uses Arnold's Cat Map for permutation and a Two-Dimensional Logistic Sine Coupling Map for diffusion. The study investigates whether structural weaknesses exist when the scheme is evaluated under a chosen-plaintext attack model. By exploiting the plaintext-independence of the keystream and the separability of permutation and diffusion stages, the work demonstrates that it is possible to recover the equivalent keystream and permutation mapping using a limited number of queries. This enables exact decryption of ciphertexts without access to the secret key, revealing critical design flaws in the encryption scheme.

\section{Overview}

\subsection{Log-Based Anomaly Detection}

The automated identification of anomalous patterns in log message sequences generated by an operational system is known as log-based anomaly detection. Semi-structured log messages combine variable parameters like block identifiers, IP addresses, error codes, and hardware locations with fixed template components. A standard HDFS message, for instance, "Receiving block blk12345 src: /10.0.2.15:50010 dest: /10.0.2.16:50010," has variable parameters that vary with each event but adheres to a consistent structural template. The fundamental challenge presented by this semi-structured nature is that anomaly detection techniques need to be sensitive to the minute patterns that differentiate abnormal behavior from typical operation while also being resilient to structural variation across logging frameworks.

The problem was approached as a two-stage pipeline in the early log anomaly detection methods. After separating constant text from variable parameters, log parsers like SLCT, PLoM~\cite{makanju2009clustering}, and Drain~\cite{he2017drain} extracted structured templates. These templates were then subjected to anomaly detection using methods like PCA on event count vectors~\cite{xu2009detecting} and invariant mining~\cite{lou2010mining}. This paradigm was conceptually simple, but it suffered from a crucial reliance on parsing accuracy. Comprehensive benchmarks by Zhu et al.~\cite{zhu2019tools} revealed that on complex real-world logs, parser error rates ranged from 14 to 89\%, with cascading effects reducing detection performance by as much as 50\%.

Advanced representation learning and sequence modeling for log anomaly detection using LSTMs, attention mechanisms, and semantic embeddings are used in deep learning techniques like DeepLog~\cite{du2017deeplog}, LogAnomaly~\cite{meng2019loganomaly}, and LogRobust~\cite{zhang2019robust}. Although these techniques greatly outperformed classical methods, they still relied on log parsing, which made them vulnerable to changing log formats and out-of-vocabulary templates.

NeuralLog~\cite{neurallog2021} introduced the parsing-free paradigm, proving that BERT~\cite{devlin2018bert} can encode raw log messages directly without requiring a parsing step in between. On benchmark datasets where parsing-based methods failed to surpass 0.60, NeuralLog achieved F1-scores above 0.95 by using WordPiece tokenization~\cite{schuster2012japanese} to handle out-of-vocabulary terms and averaging contextual token embeddings to generate fixed-dimensional message representations. This demonstrated that pre-trained language models can naturally capture the semantics of raw log text and that parsing is not a necessary precondition for successful log anomaly detection.

However, NeuralLog and current parsing-free methods have two key limitations. First, they treat all tokens within a log message the same way by simply averaging their contextual embeddings. This approach mixes tokens with different functional roles, such as states, entities, actions, quantities, and locations, into one undifferentiated vector. Second, they discard numeric tokens, believing they add noise. This oversight removes important anomaly indicators, such as HTTP error codes, memory usage values, and hardware addresses. To overcome these issues, we propose a role-aware representation learning framework that automatically identifies token roles from sequence-level anomaly labels without manual annotation or log parsing. The role induction network assigns a soft probability distribution over five role categories for each token. Role-aware pooling then creates structured message representations that maintain both semantic and structural information.

\subsection{Autoencoder-Based Attacks on Secret Image Sharing Schemes}

Secret image sharing schemes distribute a confidential image $S$ among $n$ participants by creating shares $\{s_1, s_2, ..., s_n\}$. To reconstruct $S$, a qualified subset of shares is required. XOR-based schemes achieve this with perfect secrecy; any strict subset of shares is statistically independent of the secret. They also avoid pixel expansion, making them efficient for practical use. However, this information-theoretic guarantee assumes that adversaries are computationally unbounded and have no prior knowledge of natural image statistics. Modern deep learning models can exploit these statistical patterns. This raises the question whether a learning-based adversary can reconstruct a meaningful approximation of $S$ from a limited subset of its XOR shares.

This project examines this threat by proposing an autoencoder-based reconstruction attack targeting XOR-generated secret shares. An encoder-decoder network with skip connections is trained to recover the secret image from a consecutive subset of $k < n$ shares, without needing any cryptographic key material or the remaining shares. The attack uses learned priors about natural image structures such as edges, textures, and semantic content to compensate for missing cryptographic information. The framework is evaluated across various benchmark image datasets using PSNR and SSIM as metrics for fidelity.

\subsection{Cryptanalysis of Medical Image Encryption Schemes}

Chaos-based image encryption schemes are widely used for securing medical images due to their ability to achieve high confusion and diffusion with low computational cost. These schemes typically follow a permutation-diffusion architecture, where pixel positions are first scrambled and then pixel values are modified using a pseudorandom keystream generated from chaotic systems. Despite strong statistical performance, many such schemes fail to provide resistance against stronger adversarial models.

This project performs a detailed cryptanalysis of a medical image encryption scheme that combines Arnold's Cat Map and a Two-Dimensional Logistic Sine Coupling Map. The analysis is conducted under the chosen-plaintext attack model, where the attacker has access to an encryption oracle. It is shown that the keystream used in the diffusion stage is independent of the plaintext, making it recoverable using a single carefully chosen input. Additionally, the permutation process can be reconstructed using structured probe images. By combining these two components, the complete encryption process can be inverted without knowledge of the secret key.

The results demonstrate that the scheme is fundamentally insecure despite having a large key space and passing standard statistical tests. This highlights the importance of designing encryption algorithms where keystream generation depends on the plaintext and where permutation and diffusion are tightly coupled. The study provides insights into common weaknesses in chaos-based encryption and emphasizes the need for rigorous cryptanalysis in evaluating security.